% Statement and cited-input packet. Target proof is separately embargoed.
% Frozen II source SHA256: 74497bfe327fae25c9cdcb17334b3bf42a9793a9df7c0a4404d97568f2f760f0
% Target owner: cdp:st:D.2.4; tag: 01GB

\begin{proposition}[Proved sign in the full normal-crossings class]\hypertarget{tag-01GB}{}\label{cdp:st:D.2.4}



For the threefolds and maps in D.2 whose reduced fibre supports
all have normal crossings, the Euler number is
\[
c_3(Y)=0,\qquad b_3(Y)=2.
\]
All multiplicities, including unequal multiplicities on a
reducible support, are allowed. In this class there are no
positive-Euler-number examples with trivial monodromy.
\end{proposition}

% Cited input cdp:st:C.2.NC.3; statement or labelled formula only.
\begin{proposition}[Uniform multiple normal-crossings fibres]\hypertarget{tag-01CA}{}\label{cdp:st:C.2.NC.3}


Let $g:Y\to C$ be a proper holomorphic map from a connected smooth compact complex threefold with $H^2(Y,\Q)=0$ to a smooth compact curve, with connected fibres and smooth general fibre of Euler characteristic zero. A fibre $mS$ with reduced normal-crossings support has
\[
 e(S)=\tau(S)\ge0.
 \tag{C.NC.8}\hypertarget{tag-00S7}{}\label{cdp:eq:C.NC.8}
\]
The normalized base change $t=s^m$ is finite étale of degree $m$ over a neighbourhood of $S$ and has reduced normal-crossings central fibre $S'$. Its possibly disconnected components do not affect the comparisons
$e(S')=m e(S)$ and $\tau(S')=m\tau(S)$.
\end{proposition}

% Cited input cdp:st:C.2.NC.5; statement or labelled formula only.
\begin{proposition}[Global positivity and the all-twist normalization obstruction in the normal-crossings class]\hypertarget{tag-01CC}{}\label{cdp:st:C.2.NC.5}


Let $g:Y\to C$ have the hypotheses of \ref{cdp:st:C.2.NC.3}, and assume every fibre reduction has normal crossings; smooth reductions are included. Then
\[
 c_3(Y)=e(Y)=\sum_{c\in C}\tau((Y_c)_{\rm red})=2r\ge0,
 \tag{C.NC.10}\hypertarget{tag-00S1}{}\label{cdp:eq:C.NC.10}
\]
where $r$ is the number of fibre supports with rational normalization. Every such support is irreducible, nonnormal, and has exactly two triple points. Equality $c_3(Y)=0$ holds exactly when every fibre support is free of triple points. If also $b_1(Y)=0$, then $r\in\{0,1\}$ and $b_3(Y)=2-2r$.

In the positive case there is an irreducible nonnormal reduced support $S$ with rational normalization $\eta:T\to S$ such that, for every $L\in\Pic(Y)$,
\[
 H^2\!\left(T,K_T\otimes\eta^*((K_Y^{-1}\otimes L)|_S)\right)\simeq\C,
 \qquad
 H^0\!\left(S,\Omega_Y^1|_S\otimes(K_Y\otimes L^{-1})|_S\right)\simeq\C.
 \tag{C.NC.11}\hypertarget{tag-00S2}{}\label{cdp:eq:C.NC.11}
\]
The scheme-theoretic fibre may be a multiple of $S$; this conclusion does not assert that its multiplicity is one. By \ref{cdp:st:C.3.5}, all these conclusions apply to every fibration under the CDP hypotheses whose fibre reductions have normal crossings, for every smooth-fibre type in \ref{cdp:st:C.3.1}.
\end{proposition}

% Cited input cdp:st:D.2.3; statement or labelled formula only.
\begin{proposition}[Proved logarithmic period obstruction]\hypertarget{tag-01GA}{}\label{cdp:st:D.2.3}



Let $h:U\to\Delta$ be a proper holomorphic map from a smooth
complex threefold, with connected fibres, smooth two-torus fibres
over $\Delta^*$, and central fibre $F=\sum_i m_iD_i$ with
normal-crossings reduced support $D=\bigcup_iD_i$.
If every normalization component of $D$ has geometric genus zero,
the monodromy on $H^2(U_t;\C)$ is nontrivial. In particular the
monodromy on $H^1(U_t;\Q)$ is nontrivial. No K\"ahler assumption
on $U$ or $D$ is made, and the positive integers $m_i$ are arbitrary.
\end{proposition}

% Cited input cdp:st:D.2.7; statement or labelled formula only.
\begin{proposition}[Proved existence]\hypertarget{tag-01GE}{}\label{cdp:st:D.2.7}



There exists a smooth compact connected complex threefold $Y$
with a holomorphic map $f:Y\to\PP^1$ with connected fibres such
that
\[
 H_1(Y;\Z)=H_2(Y;\Z)=0,\qquad
 a(Y)=1,\qquad c_3(Y)=0,\qquad b_3(Y)=2.
\tag{D.21}\hypertarget{tag-0198}{}\label{cdp:eq:D.21}
\]
Every smooth fibre is the fixed abelian surface
$A=E\times E$, where $E=\C/(\Z+\mathrm i\Z)$, and the marked
integral monodromy on $H_1(A;\Z)$ is the identity. There are
exactly three singular fibres. Each is reduced with simple
normal crossings, and consists of a primary Kodaira surface
and a primary Hopf surface meeting along one smooth elliptic
curve. Consequently the trivial-monodromy class in D.2 is
nonempty, and equality in the unrestricted sign theorem
\ref{cdp:st:D.2.5} is attained.
\end{proposition}

% Cited input cdp:st:C.3.1; statement or labelled formula only.
\begin{proposition}\hypertarget{tag-01F4}{}\label{cdp:st:C.3.1}


Let $f:Y\to C$ be a holomorphic algebraic reduction of a connected
compact complex threefold with $b_1(Y)=b_2(Y)=0$ and $a(Y)=1$.
The base is $C\simeq\PP^1$. Every smooth fibre $F$ has
\[
 c_1(K_F)=0\text{ in }H^2(F;\Q),\qquad
 c_2(F)=0\text{ in }H^4(F;\Z),\qquad
 \chi(F,\OO_F)=0,
\]
and contains no exceptional curve of the first kind. A very general
smooth fibre has $\kappa(F)\leq0$. The very general fibre's type belongs to the following
list: a complex two-torus, a bielliptic surface, a primary or secondary
Kodaira surface, a primary or secondary Hopf surface, or an Inoue
surface. Every listed type has Euler characteristic zero. This is a
preliminary necessary list. Statement \ref{cdp:st:C.3.7} below excludes the
bielliptic entry, and \ref{cdp:st:C.3.8} makes the very general Hopf fibre have
algebraic dimension zero. \ref{cdp:st:C.3.9} and \ref{cdp:st:C.3.10} exclude both Kodaira
entries. Statement \ref{cdp:st:C.3.6} realizes both primary and secondary Hopf
entries, and \ref{cdp:st:D.2.7} realizes a fixed abelian surface as general
fibre. Statement \ref{cdp:st:C.3.12} excludes the Inoue entry. Thus the
final very general list consists exactly of two-tori and
primary or secondary Hopf surfaces of algebraic dimension zero.
\end{proposition}

% Cited input cdp:st:C.3.5; statement or labelled formula only.
\begin{proposition}\hypertarget{tag-01FB}{}\label{cdp:st:C.3.5}


For every threefold and fibration in \ref{cdp:st:C.3.1}, including every type on
its necessary list, the Euler-number identity is
\[
\int_Y c_3(Y)=e(Y)=\sum_{c\in\Delta}e\bigl((Y_c)_{\rm red}\bigr),
\]
where $\Delta$ contains the finitely many critical values.
The identity alone assigns no sign to the summands.
\end{proposition}

% Cited input cdp:st:D.2.5; statement or labelled formula only.
\begin{proposition}[Proved unrestricted sign under trivial monodromy]\hypertarget{tag-01GC}{}\label{cdp:st:D.2.5}



Let $f:Y\to\PP^1$ be a proper surjective holomorphic map with
connected fibres from a compact connected complex threefold with
$b_1(Y)=b_2(Y)=0$. Suppose its smooth general fibre is a complex
two-torus and its monodromy representation on $H_1(F,\Q)$ is
trivial. Then
\[
 c_3(Y)\le0,\qquad b_3(Y)\ge2.
 \tag{D.S.1}\hypertarget{tag-019O}{}\label{cdp:eq:D.S.1}
\]
The fibres may have arbitrary singularities and arbitrary
multiplicities. In particular this proves the sign assertion of
D.2 for a holomorphic algebraic reduction; no existence assertion
is required for the conclusion.
\end{proposition}

% Cited input cdp:st:C.3.10; statement or labelled formula only.
\begin{proposition}[Proved exclusion of primary Kodaira general fibres]\hypertarget{tag-01F5}{}\label{cdp:st:C.3.10}



A compact complex threefold with $b_1(Y)=b_2(Y)=0$ and
$a(Y)=1$ cannot have a holomorphic algebraic reduction
$f:Y\to\PP^1$ whose general fibre is a primary Kodaira
surface. Together with \ref{cdp:st:C.3.9} this excludes both Kodaira
types using only rational Betti-number vanishing.
\end{proposition}

% Cited input cdp:st:C.3.12; statement or labelled formula only.
\begin{proposition}[Proved exclusion of Inoue general fibres]\hypertarget{tag-01F7}{}\label{cdp:st:C.3.12}



Under the hypotheses of \ref{cdp:st:C.3.1}, the very general smooth fibre is
not an Inoue surface. More precisely, a proper holomorphic map from
a compact connected complex manifold to a compact Riemann surface,
whose fibres over the complement of a finite set are Inoue surfaces,
has a class in $H^1(Y;\mathbb R)$ restricting nontrivially to each
such fibre. Thus these families cannot have $b_1(Y)=0$.
This assertion permits arbitrary multiplicities and arbitrary
singularities of the omitted fibres.
\end{proposition}

% Cited input cdp:st:C.3.6; statement or labelled formula only.
\begin{proposition}[Proved actual occurrence of both Hopf types]\hypertarget{tag-01FC}{}\label{cdp:st:C.3.6}



There are compact complex threefolds with $b_1=b_2=0$ and $a=1$
whose holomorphic algebraic reduction is a smooth bundle over
$\PP^1$ with primary Hopf fibre. There are also such threefolds
with secondary Hopf fibre. In these examples every fibre has Euler
number zero and $c_3=0$.
\end{proposition}

% Cited input cdp:st:C.3.7; statement or labelled formula only.
\begin{proposition}[Proved exclusion of bielliptic fibres]\hypertarget{tag-01FD}{}\label{cdp:st:C.3.7}



A compact complex threefold $Y$ with $b_1(Y)=b_2(Y)=0$ has no
proper connected-fibre holomorphic map to a compact curve whose
general fibre is bielliptic.
\end{proposition}

% Cited input cdp:st:C.3.8; statement or labelled formula only.
\begin{proposition}[Proved exclusion of elliptic Hopf general fibres]\hypertarget{tag-01FE}{}\label{cdp:st:C.3.8}



The very general Hopf fibre of the algebraic reduction in \ref{cdp:st:C.3.1}
has algebraic dimension zero. In particular neither primary nor
secondary Hopf surfaces of algebraic dimension one can be its
general fibre. This exclusion uses no restriction-surjectivity
assertion from CDP98 or CDP20.
\end{proposition}

% Cited input cdp:st:C.3.9; statement or labelled formula only.
\begin{proposition}[Proved exclusion of secondary Kodaira general fibres]\hypertarget{tag-01FF}{}\label{cdp:st:C.3.9}



A compact complex threefold with $b_1(Y)=b_2(Y)=0$ and
$a(Y)=1$ cannot have a holomorphic algebraic reduction
$f:Y\to\PP^1$ whose general fibre is a secondary Kodaira
surface. Rational Betti-number vanishing suffices; no integral
torsion-freeness of the fibre is assumed.
\end{proposition}

